\begin{abstract}
Deep-learning drowsiness detectors are hard to trust in a vehicle because their decisions are opaque, and interpretable and explainable methods are offered as the remedy. We review 45 papers across electroencephalography (EEG), cardiac signals, camera-based vision, vehicle signals and multimodal fusion, classify them by how the explanation is produced and what it refers to, and appraise 41 detection studies on four checks: agreement with known physiology, a faithfulness test, testing on unseen drivers, and public code. Thirty-two studies report explanations at least partly consistent with known physiology, but only two tested faithfulness quantitatively, 28 did not evaluate the explanation, 15 clearly tested on unseen drivers, and five share code. Where explanations were tested they could disagree: attention and band ablation conflict in one EEG model, and four explainers rank differently on one alertness model. Several labels derive from model inputs or recording sessions, and median reported accuracy is 97.1\% under random or within-subject splits against 77.2\% on unseen drivers. A case study on a published interpretable EEG model shows that it outputs a constant class when normalisation cannot see the whole test subject, that changing only those statistics alters up to 24.6\% of decisions, that attribution methods strongest on deletion are weakest on a randomisation sanity check, and that the reading of alert decisions as artefacts holds weakly on the 11-subject subset but more clearly on all 27 subjects. We close with a reporting checklist. The papers came from targeted searches, so counts describe the examined set, not the field.
\end{abstract}

\begin{IEEEkeywords}
Drowsiness detection, driver monitoring, explainable artificial intelligence, interpretability, electroencephalography, heart rate variability, faithfulness, cross-subject validation.
\end{IEEEkeywords}
